\chapter{Quadratic logistic regression analysis}

In this section, we move beyond linear decision boundaries by applying a quadratic feature expansion to the project dataset. By mapping the original features into a higher-dimensional space that includes all pairwise products of the features ($\phi(x) = [\text{vec}(xx^T), x]^T$), we allow the Logistic Regression model to learn complex, non-linear boundaries in the original feature space.

\section{Experimental results}

The model was trained on the full expanded dataset using a standard (non-prior-weighted) objective function. We evaluated the performance across the same 13 logarithmically spaced $\lambda$ values for the primary application $\pi_T = 0.1$. 

The numerical results are summarized in Table \ref{tab:quad_lr_results} and the visual trends are plotted in Figure \ref{fig:lr_quad_plot}.

\begin{table}[h]
\centering
\caption{Quadratic logistic regression results ($\pi_T = 0.1$)}
\label{tab:quad_lr_results}
\begin{tabular}{@{}ccc@{}}
\toprule
\textbf{Regularization Coefficient ($\lambda$)} & \textbf{minDCF} & \textbf{actDCF} \\ \midrule
$1.0000 \times 10^{-4}$ & 0.2602 & 0.2768 \\
$3.1623 \times 10^{-4}$ & 0.2612 & \textbf{0.2656} \\
$1.0000 \times 10^{-3}$ & 0.2587 & 0.2765 \\
$3.1623 \times 10^{-3}$ & 0.2527 & 0.2771 \\
$1.0000 \times 10^{-2}$ & 0.2487 & 0.3464 \\
$3.1623 \times 10^{-2}$ & \textbf{0.2436} & 0.4972 \\
$1.0000 \times 10^{-1}$ & 0.2466 & 0.7520 \\
$3.1623 \times 10^{-1}$ & 0.2629 & 0.9623 \\
$1.0000 \times 10^{0}$  & 0.2840 & 1.0000 \\
$3.1623 \times 10^{0}$  & 0.3090 & 1.0000 \\
$1.0000 \times 10^{1}$  & 0.3243 & 1.0000 \\
$3.1623 \times 10^{1}$  & 0.3263 & 1.0000 \\
$1.0000 \times 10^{2}$  & 0.3263 & 1.0000 \\
\bottomrule
\end{tabular}
\end{table}

\begin{figure}[h]
    \centering
    \includegraphics[width=0.85\textwidth]{LR_Quadratic_DCF_vs_Lambda.png}
    \caption{Actual DCF and Minimum DCF as a function of $\lambda$ for the Quadratic Logistic Regression model ($\pi_T = 0.1$).}
    \label{fig:lr_quad_plot}
\end{figure}

The optimal discriminative capability is reached at an intermediate regularization of $\lambda = 3.16 \times 10^{-2}$, yielding an outstanding minDCF of \textbf{0.2436}. However, the best calibration occurs at a much lower regularization ($\lambda = 3.16 \times 10^{-4}$), where actDCF drops to an excellent \textbf{0.2656}.

\section{Analysis and discussion}

\subsection*{What do you observe?}

The most striking observation is the massive, systemic improvement in discriminatory power compared to all previously evaluated linear models. The minimum DCF drops precipitously from approximately $0.361$ in the linear case down to a best value of $0.2436$ in this quadratic case. This mathematically proves that the relationship between the features and the classes is inherently non-linear, and the quadratic expansion successfully maps this complexity.

Furthermore, we observe that the minDCF curve (blue line in Figure \ref{fig:lr_quad_plot}) is no longer flat as it was in the linear full-dataset analysis. Instead, it exhibits a distinct "U-shape," slowly improving as $\lambda$ increases toward $10^{-2}$ before deteriorating as the penalty becomes too aggressive.

The actDCF (red line) behaves exceptionally well for small values of $\lambda \le 10^{-3}$, remaining tightly coupled to the minDCF curve and showing excellent initial calibration (e.g., $0.2656$ actDCF vs $0.2612$ minDCF). However, it degrades sharply as regularization scales up, hitting the $1.0000$ ceiling exactly like the linear models.

\subsection*{In this case is regularization effective?}

Yes, in the quadratic case, explicit regularization is highly effective and absolutely necessary for optimizing discriminative power. 

In the standard linear analysis on the full dataset, the number of parameters was tiny (6 features + 1 bias = 7 weights). With thousands of training samples, the model was immune to variance, making $L_2$ regularization largely irrelevant for the minDCF. However, the quadratic expansion drastically increases the dimensionality of the feature space. Computing the outer product $xx^T$ inflates the parameter count substantially, providing the model with a massively expanded capacity.

With this high capacity comes a severe vulnerability to overfitting. We can see regularization counteracting this by observing the minDCF: it improves from $0.2602$ at $\lambda=10^{-4}$ down to the global optimum of $0.2436$ at $\lambda=3.16 \times 10^{-2}$. This proves that penalizing the expanded weight vector prevents the model from mapping the random noise generated by the squared feature terms, locking in a more robust geometric boundary.

\subsection*{How does it affect the two metrics?}

The regularization coefficient $\lambda$ controls the bias-variance tradeoff, affecting the metrics as follows:

\begin{itemize}
    \item \textbf{Effect on minDCF:} For very low $\lambda$, the high-capacity model slightly overfits the training set (higher minDCF due to variance). As $\lambda$ increases toward $3.16 \times 10^{-2}$, the minDCF reaches its absolute optimum because the $L_2$ penalty forces the model to ignore spurious quadratic interactions while preserving the core non-linear class signals. Beyond this point, the model begins to heavily underfit (higher minDCF due to bias), as the massive penalty prevents the weights from drawing any meaningful boundaries.
    
    \item \textbf{Effect on actDCF:} For low $\lambda$, calibration is remarkably accurate because the expanded quadratic model possesses the structural complexity required to accurately model the true log-posterior probabilities without needing severe weight suppression. However, as $\lambda$ increases beyond the optimal point, the actDCF explodes. Just as in the linear variants, heavy regularization forcefully shrinks the complex weight vector toward zero, artificially compressing the output scores toward the bias term. These compressed scores lose their alignment with the theoretical Bayes threshold ($t \approx 2.19$ for $\pi_T = 0.1$), resulting in a massive increase in the empirical cost.
\end{itemize}

\section{Theoretical motivation}

The overwhelming superiority of the quadratic Logistic Regression model ($0.2436$ minDCF) stems from its ability to model covariance differences between classes using a purely discriminative optimization approach. 

A standard linear logistic regression is mathematically analogous to a Tied Covariance Gaussian model, in that both are strictly constrained to drawing a flat, linear hyperplane. By passing the outer product $xx^T$ into the model, the quadratic expansion allows the Logistic Regression algorithm to natively compute boundary curvatures, effectively simulating the behavior of a fully unconstrained Multivariate Gaussian (MVG) model. 

In fact, the Quadratic LR outperforms our baseline Generative MVG (which scored a $0.2629$ minDCF). This happens because the Generative MVG strictly assumes the underlying data follows a perfect Normal distribution—an assumption we know is violated by the heavy skewness of Features 5 and 6 in our dataset. The Quadratic Logistic Regression algorithm learns the non-linear boundaries directly from the error gradients, entirely bypassing the rigid, flawed Gaussian density assumption, thereby extracting the absolute maximum discriminative utility from the feature space.